% Figure: the Hyatt Regency hanger-rod connection. Original single-rod
% load path (left) versus the as-built two-rod load path (right) that
% doubled the upper connection load.
\begin{figure}[htbp]
\centering
\begin{tikzpicture}[
  rod/.style={draw=enggray, very thick},
  beam/.style={draw=engblack, fill=engblue!10, thick},
  load/.style={draw=engred, -{Stealth}, very thick},
  scale=1.0,
]
% ===== Left: original design, single continuous rod =====
\begin{scope}[xshift=0cm]
  % ceiling
  \draw[draw=engblack, very thick] (-0.4,5.4) -- (3.0,5.4);
  \foreach \x in {-0.2,0.2,...,2.8} {\draw[draw=enggray, thin] (\x,5.4) -- (\x-0.2,5.65);}
  % single rod from ceiling all the way down
  \draw[rod] (1.3,5.4) -- (1.3,0.6);
  % upper (4th-floor) box beam at y=3.6
  \draw[beam] (0.2,3.4) rectangle (2.4,3.85);
  \node[font=\scriptsize] at (1.3,3.62) {4th-floor beam};
  % nut at upper beam
  \draw[draw=engblack, fill=engamber!50] (1.15,3.4) rectangle (1.45,3.55);
  % lower (2nd-floor) box beam at y=1.0
  \draw[beam] (0.2,0.8) rectangle (2.4,1.25);
  \node[font=\scriptsize] at (1.3,1.02) {2nd-floor beam};
  \draw[draw=engblack, fill=engamber!50] (1.15,0.6) rectangle (1.45,0.75);
  % loads
  \draw[load] (0.5,3.85) -- (0.5,4.55) node[anchor=south, font=\scriptsize, engred] {$W_4$};
  \draw[load] (0.5,1.25) -- (0.5,1.95) node[anchor=south, font=\scriptsize, engred] {$W_2$};
  % connection load callout
  \node[font=\scriptsize, anchor=west, engblue] at (2.6,3.6) {carries $W_4$};
  \node[font=\small, anchor=north] at (1.3,0.0) {(a) original: single rod};
\end{scope}
% ===== Right: as-built, two rods, doubled connection load =====
\begin{scope}[xshift=7.4cm]
  \draw[draw=engblack, very thick] (-0.4,5.4) -- (3.0,5.4);
  \foreach \x in {-0.2,0.2,...,2.8} {\draw[draw=enggray, thin] (\x,5.4) -- (\x-0.2,5.65);}
  % upper rod: ceiling to 4th-floor beam
  \draw[rod] (1.3,5.4) -- (1.3,3.85);
  % lower rod: 4th-floor beam to 2nd-floor beam (OFFSET to show it is separate)
  \draw[rod] (1.7,3.4) -- (1.7,1.25);
  % upper beam
  \draw[beam] (0.2,3.4) rectangle (2.4,3.85);
  \node[font=\scriptsize] at (0.9,3.62) {4th-floor};
  \draw[draw=engblack, fill=engamber!50] (1.15,3.4) rectangle (1.45,3.55);
  \draw[draw=engblack, fill=engamber!50] (1.55,3.4) rectangle (1.85,3.55);
  % lower beam
  \draw[beam] (0.2,0.8) rectangle (2.4,1.25);
  \node[font=\scriptsize] at (0.9,1.02) {2nd-floor};
  \draw[draw=engblack, fill=engamber!50] (1.55,1.25) rectangle (1.85,1.4);
  % loads
  \draw[load] (0.5,3.85) -- (0.5,4.55) node[anchor=south, font=\scriptsize, engred] {$W_4$};
  \draw[load] (2.2,1.25) -- (2.2,1.95) node[anchor=south, font=\scriptsize, engred] {$W_2$};
  % connection load callout: now W4 + W2
  \node[font=\scriptsize, anchor=west, engred] at (2.6,3.6) {carries $W_4 + W_2$};
  \node[font=\small, anchor=north] at (1.3,0.0) {(b) as-built: two rods};
\end{scope}
\end{tikzpicture}
\caption[Hyatt Regency: the load path that the redrawn FBD would have
caught]{The hanger-rod connection at the heart of the 1981 Hyatt
Regency walkway collapse \cite{acc:nbs-hyatt-1982}. (a) In the original
design a single continuous rod runs from the ceiling through both
walkway box-beams; the upper connection carries only the fourth-floor
walkway load $W_4$, because the second-floor load $W_2$ reaches the
ceiling through the rod below the upper nut. (b) The as-built design,
changed during fabrication to ease assembly, used two rod segments: the
lower walkway now hangs from the upper walkway's box-beam, so the upper
connection carries $W_4 + W_2$, roughly double the design intent. The
change looks local on a drawing. The free-body diagram of the upper
connection, redrawn for the as-built geometry, makes the doubled load
visible at once. That diagram was never drawn.}
\label{fig:vol03ch01:hyatt-loadpath}
\end{figure}
